\chapter{The Spectral Theorem and Rayleigh Quotients}
\label{ch:spectral-theorem-rayleigh}

\begin{draftbox}
This chapter is scaffolded. It will absorb the spectral theorem proof,
invariant subspaces, and Rayleigh quotient lab material.
\end{draftbox}

\section{Guiding Question}

Why are symmetric matrices the best-behaved matrices in real linear algebra?

\section{Planned Core Ideas}

\begin{itemize}
  \item symmetric matrices preserve orthogonal complements of invariant
    subspaces;
  \item the spectral theorem gives an orthonormal eigenbasis;
  \item PSD matrices have nonnegative eigenvalues;
  \item the Rayleigh quotient turns eigenvalues into optimization problems;
  \item numerical eigenvalue methods exploit this structure.
\end{itemize}

\section{Planned Python Thread}

Implement simple Rayleigh quotient experiments and compare with
\texttt{np.linalg.eigh}.

\section*{Chapter Summary}
\addcontentsline{toc}{section}{Chapter Summary}

To be completed in the chapter-shaping pass.

\section*{Exercises}
\addcontentsline{toc}{section}{Exercises}
\subsection*{A. Theory}
\begin{exercise}\exerciseid{15.A1} State the spectral theorem for real symmetric matrices and explain why orthogonality is part of the conclusion. \end{exercise}
\subsection*{B. By Hand}
\begin{exercise}\exerciseid{15.B1} Compute the Rayleigh quotient of $\mat{2&0\\0&5}$ at $(1,1)$ and at $(0,1)$. \end{exercise}
\subsection*{C. Python / NumPy}
\begin{exercise}\exerciseid{15.C1} Sample random unit vectors and evaluate the Rayleigh quotient for a symmetric $2\times 2$ matrix. \end{exercise}
\subsection*{D. Challenge}
\begin{exercise}\exerciseid{15.D1} Explain why the maximum Rayleigh quotient of a symmetric matrix should be an eigenvalue. \end{exercise}

